\chapter{Family--defect projector supercurvature factorization}

Предыдущий gate дал coefficient-free zero locus, но записал radial и axis
условия разными слагаемыми. Проверим более сильную возможность: получить
condensate и тетраэдрическую family axis из одной quadratic curvature.

\section{Перепараметризация winding branch}

Положим \(y=\nu x\). Тогда прежнее условие
\(I_3(x)=\nu/\sqrt3\) принимает вид
\begin{equation}
 I_3(y)=\frac1{\sqrt3}.
\end{equation}
Ось \(y\) всегда лежит в одной четырёхточечной тетраэдрической орбите,
а winding задаёт направление поворота
\begin{equation}
 R_{y,\nu}=R_y\!\left(\nu\frac{2\pi}{3}\right).
\end{equation}

\section{Одна shifted-square curvature}

Для coordinate rank-one projector \(P_a\) на four-state affine menu
введём
\begin{equation}
 H_a=\frac2{\sqrt3}\left(P_a-\frac14I_4\right)
\end{equation}
и curvature polynomial
\begin{equation}
 Q(H)=H^2-\frac1{\sqrt3}H-\frac14I_4.
 \label{eq:projector-supercurvature}
\end{equation}
Он факторизуется как
\begin{equation}
 Q(H)=\left(H-\frac1{2\sqrt3}I_4\right)^2-\frac13I_4.
\end{equation}
При \(J=\sqrt3H-I_4/2\) equation \(Q(H)=0\) эквивалентна
\(J^2=I_4\). Tracelessness даёт \(\Tr J=-2\), поэтому
\begin{equation}
 P=\frac{I_4+J}{2}
\end{equation}
автоматически является rank-one projector.

\begin{theorem}[Projector factorization]
На diagonal traceless \(4\times4\) matrices equation \(Q(H)=0\)
имеет ровно четыре решения
\begin{equation}
 H_a=\frac2{\sqrt3}\left(P_a-\frac14I_4\right),
 \qquad a=1,\ldots,4.
\end{equation}
Они удовлетворяют
\begin{equation}
 \Tr H_a^2=1,
 \qquad
 \Tr H_a^3=\frac1{\sqrt3}.
\end{equation}
Следовательно, одна curvature equation фиксирует radial norm и одну из
четырёх тетраэдрических axes.
\end{theorem}

Каждый eigenvalue является корнем
\(h^2-h/\sqrt3-1/4=0\). Поэтому возможны только
\(h=\sqrt3/2\) и \(h=-1/(2\sqrt3)\); нулевой trace оставляет один
eigenvalue первого типа и три второго.

Для \(V_Q=\Tr Q(H)^2\) все три Hessian eigenvalues на traceless diagonal
tangent space равны
\begin{equation}
 \lambda_1=\lambda_2=\lambda_3=\frac83.
\end{equation}
Radial и две axis directions стабилизируются одной trace-нормой.

\section{От projector к oriented three-cycle}

Projector \(P_a\) выбирает fixed vertex. На оставшихся трёх вершинах
существуют ровно два inverse three-cycles; их различает winding
\(\nu=\pm1\). Обозначим их через \(C_{a,\nu}\). Для любого
\(g\in S_4\) выполняется twisted covariance law
\begin{equation}
 gC_{a,\nu}g^{-1}
 =C_{g(a),\operatorname{sgn}(g)\nu}.
 \label{eq:twisted-cycle-covariance}
\end{equation}
Все \(24\times4\times2=192\) cases проверены машинно; failures нет.

На standard triplet каждая permutation имеет
\begin{equation}
 \Tr R_{a,\nu}=0,
 \qquad
 \det R_{a,\nu}=1,
\end{equation}
и равна
\begin{equation}
 R_{a,\nu}
 =\exp\!\left(\nu\frac{2\pi}{3}[n_a]_\times\right).
\end{equation}
Соответствующий Majorana generator имеет rank два и nullity один.

\begin{proposition}[Projector--cycle reconstruction]
Четыре curvature-selected projectors и два winding orientations
реконструируют все восемь \(S_4\) three-cycles. Реконструкция ковариантна,
если winding является orientation-line variable по
rule~\eqref{eq:twisted-cycle-covariance}.
\end{proposition}

\section{Статус parent action}

Теперь radial equation и tetrahedral axis не являются независимыми
potentials: обе следуют из shifted involution. Связь с four-state carrier
также прямая, поскольку solutions являются его primitive projectors.
Superconnection architecture допускает odd endomorphism и curvature square
для position-dependent mass, см. \path{arXiv:2106.01591}; specific
factorization здесь получена из project affine menu.

Но equation~\eqref{eq:projector-supercurvature} пока не является
вычисленным блоком полного \(\mathbb A^2\). Нужно показать, что:
\begin{enumerate}
 \item \(Q(H)\) возникает как projected even curvature;
 \item physical winding является orientation-line variable;
 \item \(C_{a,\nu}\) возникает как actual holonomy, а не как
 post-processing projector;
 \item та же connection ограничивается на Majorana kernel как
 \(K_{a,\nu}=\nu(2\pi/3)[n_a]_\times\).
\end{enumerate}

\begin{nogocriterion}[Holonomy realization gate]
Projector factorization не является полным boundary-superconnection
выводом, пока \(C_{a,\nu}\) не получена как holonomy той же connection,
чья curvature содержит \(Q(H)\).
\end{nogocriterion}

Следующий flat-holonomy gate проходит этот критерий на constrained saddle:
canonical connection, построенная из \(H_a\), uniform singlet и
orientation tensor, имеет holonomy ровно \(C_{a,\nu}\). Открытым остаётся
уже не holonomy realization, а вариационное происхождение constitutive
relation между connection и projector field.

Итак, требуется вывести уже не три locking terms, а один projected
curvature block и его oriented holonomy.

Машинный ledger сохранён в
\path{s2t_v4_family_defect_projector_supercurvature_gate_results.json};
генератор ---
\path{s2t_v4_family_defect_projector_supercurvature_gate.py}.